% ============================================================
% Section 3: Theoretical Analysis
% ============================================================

\section{Theoretical Analysis}
\label{sec:theory}

We establish that MIA advantage is upper-bounded by data-distributional
quantities computable before training.
All proofs are in Appendix~\ref{app:proofs}.

\subsection{Notation}

Let $D = \{x_1,\dots,x_n\} \subset \mathcal{X}$ be drawn i.i.d.\ from
density $p$ on $\mathcal{X} \subseteq \mathbb{R}^d$.
For any $x \in D$, define:
\begin{align}
  u(x) &= \min_{x' \in D,\, x'\neq x} \|x - x'\|_2
           \quad \text{(sample uniqueness / nearest-neighbour distance)},
           \label{eq:uniqueness}\\
  \rho(x) &= \frac{k}{ n \cdot \mathrm{Vol}(B(x,\, r_k(x)))}
           \quad \text{(local density, $k$-NN estimate)},
           \label{eq:density}
\end{align}
where $r_k(x)$ is the distance to the $k$-th nearest neighbour of $x$ in
$D \setminus \{x\}$ and $\mathrm{Vol}(\cdot)$ denotes Euclidean ball volume.

For the in-distribution and out-distribution of a trained model
$f_\theta$, define the score distributions:
\begin{equation}
  P_{\mathrm{in}}(s) = \Pr[\ell(f_\theta, x) \leq s \mid x \in D], \quad
  P_{\mathrm{out}}(s) = \Pr[\ell(f_\theta, x) \leq s \mid x \notin D],
\end{equation}
where $\ell(f_\theta,x)$ is the model's loss on sample $x$.

\subsection{Main Theorem}

\begin{theorem}[DPRI Upper Bound on MIA Advantage]
\label{thm:main}
Let $f_\theta$ be trained on $D$ by minimising empirical risk.
Assume the loss function $\ell$ is $L$-Lipschitz in $x$.
Then for any sample $x \in D$:
\begin{equation}
  \mathrm{Adv}(\mathcal{A}^*, f_\theta, x)
  \;\leq\;
  \sqrt{\frac{L^2 \cdot u(x)^2}{2\,\sigma^2}}
  \;\cdot\;
  \frac{1}{\rho(x)^{1/d}},
  \label{eq:bound}
\end{equation}
where $\mathcal{A}^*$ is the Bayes-optimal adversary,
$\sigma^2$ is the variance of the loss under $P_{\mathrm{out}}$,
and $d$ is the data dimensionality.

At the dataset level, the aggregate advantage satisfies:
\begin{equation}
  \mathrm{Risk}(D)
  \;\leq\;
  \frac{1}{n}\sum_{x \in D}
  \sqrt{\frac{L^2 \cdot u(x)^2}{2\,\sigma^2\,\rho(x)^{2/d}}}.
  \label{eq:risk_bound}
\end{equation}
\end{theorem}

\paragraph{Interpretation.}
The bound says that privacy risk increases when samples are
far from their neighbours ($u(x)$ large, i.e.\ isolated)
and when the local population is sparse ($\rho(x)$ small).
Both quantities are computable directly from $D$, with no model training required.
The Lipschitz constant $L$ and noise level $\sigma$ are model-family
constants that can be estimated or treated as fixed hyperparameters.

\begin{proof}[Proof sketch]
We proceed in three steps.

Step 1: KL to advantage.
By the Bretagnolle--Huber inequality, the Bayes-optimal MIA advantage is:
\begin{equation}
  \mathrm{Adv}(\mathcal{A}^*) \leq
  \sqrt{\tfrac{1}{2}\,\mathrm{KL}(P_{\mathrm{in}} \| P_{\mathrm{out}})}.
  \label{eq:bh}
\end{equation}

Step 2: KL in terms of loss gap.
Since $P_{\mathrm{in}}$ and $P_{\mathrm{out}}$ are loss score distributions,
and $\ell$ is $L$-Lipschitz, the expected loss gap between a training
sample $x$ and the nearest out-of-sample point $x'$ satisfies:
\begin{equation}
  \mathbb{E}[\ell(f_\theta, x)] - \mathbb{E}[\ell(f_\theta, x')]
  \;\leq\; L \cdot \|x - x'\|_2
  \;\leq\; L \cdot u(x).
  \label{eq:lipschitz_gap}
\end{equation}
Under a sub-Gaussian model for $P_{\mathrm{out}}$ with variance $\sigma^2$:
\begin{equation}
  \mathrm{KL}(P_{\mathrm{in}} \| P_{\mathrm{out}})
  \;\leq\;
  \frac{(L \cdot u(x))^2}{2\,\sigma^2}.
  \label{eq:kl_bound}
\end{equation}

Step 3: Density correction.
The nearest-neighbour distance $u(x)$ and local density $\rho(x)$ are
related by the $k$-NN geometry: for a sample in a region of density
$\rho$, the expected distance to the nearest neighbour scales as
$\rho(x)^{-1/d}$ (standard result, see \cite{biau2015lectures}).
Substituting into~\eqref{eq:kl_bound} and applying~\eqref{eq:bh}
yields~\eqref{eq:bound}.
Averaging over all $x \in D$ gives~\eqref{eq:risk_bound}.
\qed
\end{proof}

\subsection{Discussion of the Bound}

\paragraph{Tightness.}
The bound is tight up to constants for datasets with a uniform distribution
over a compact manifold (Proposition~\ref{prop:tight}, Appendix~\ref{app:proofs}).
For real-world datasets, the bound overestimates risk by a factor that
decreases as the data distribution becomes more regular.
Section~\ref{sec:results} shows empirically that the bound ranks datasets
in the same order as the observed AUC (Spearman $\rho > 0.85$), suggesting
it is tight enough to be predictively useful.

\paragraph{What the bound cannot capture.}
The bound depends on $L$ (Lipschitz constant of the loss) and $\sigma$
(out-distribution noise), which are model-family properties.
This means the bound explains why dataset structure is the dominant
driver of risk, while model architecture contributes a fixed multiplicative
constant, consistent with our empirical Finding~1 (Section~\ref{sec:finding1}).

\paragraph{Relationship to differential privacy.}
If the training mechanism is $(\epsilon,\delta)$-DP, then
$\mathrm{Adv}(\mathcal{A}^*) \leq e^\epsilon - 1 + \delta$~\cite{yeom2018privacy}.
Combined with Theorem~\ref{thm:main}, high-DPRI datasets require
smaller $\epsilon$ to achieve the same risk reduction, motivating our
data-adaptive DP calibration recommendation (Section~\ref{sec:finding3}).

\subsection{Corollary: Benchmark Bias is Predictable}

\begin{corollary}[Benchmark datasets have structurally elevated DPRI]
\label{cor:benchmark}
If a benchmark dataset $D_b$ was constructed by subsampling or clustering
a larger dataset to create balanced, well-separated classes, then
$\mathrm{DPRI}(D_b) > \mathrm{DPRI}(D_{\mathrm{real}})$
for a real-world dataset $D_{\mathrm{real}}$ drawn from the same domain,
in expectation over the construction procedure.
\end{corollary}

\begin{proof}[Proof sketch]
Subsampling increases $u(x)$ (nearest neighbours are removed) and
clustering increases class separation, both of which increase the RHS
of~\eqref{eq:risk_bound}.
\end{proof}

This corollary provides the theoretical underpinning for our empirical
Finding~2 (benchmark bias, Section~\ref{sec:finding2}).
